\section{Moduł BMP280}
Moduł BMP280 jest obsługiwany przez strukturę \texttt{Bmp280},
pokazaną na listingu~\ref{lst:bmp280_struct}. Zawiera ona pole
adresu urządzenia magistrali \isqrc, pole przechowujące dane
kompensacyjne pomiarów oraz parametr typu \texttt{CONFIG} zawarty
w~polu \texttt{PhantomData} o~zerowym rozmiarze w~pamięci.

\begin{lstlisting}[language=Rust, style=colouredRust, 
    caption={Struktura \texttt{Bmp280}},
    label={lst:bmp280_struct}
]
pub struct Bmp280<CONFIG>{
    address: u8,
    compensations: Option<Bmp280Compensations>,
    _config: PhantomData<CONFIG>,
}
\end{lstlisting}

\subsection*{Dane kompensacyjne}
Do poprawnej interpretacji danych modułu BMP280 niezbędne są dane kompensacyjne
zapisane przez producenta w~pamięci samego modułu.
Składają sie one z~dwunastu liczb szesnastobitowych (dwie z~nich traktowane są
jako liczby ze znakiem) zapisanych w formacie \textit{little endian}.

Przechowaniem ich po stronie programu zajmuje się struktura 
\texttt{Bmp280Compensations}.
Została ona opatrzona metodą pozwalającą na ich odczyt, przedstawioną na
listingu~\ref{lst:bmp280compensations_read}.
Metoda odczytuje 24~bajty z~pamięci modułu BMP280 i~przypisuje je do 
poszczególnych pól struktury. Nazwy pól struktury przyjęto zgodnie z~nazwami
prezentowanymi w~nocie katalogowej~\cite{bmp280_datasheet}.

\begin{lstlisting}[language=Rust, style=colouredRust, 
    caption={Odczyt danych kompensacyjnych dla modułu BMP280},
    label={lst:bmp280compensations_read}
]
pub fn read_from_i2c(address: u8, i2c: &mut I2c) 
    -> Result<Self, i2c::Error>{

    let mut buffer = [0u8; N_COMPENSATION_BYTES];

    i2c.write_read(
        address, 
        &[COMPENSATION_START_ADDRESS], 
        &mut buffer
    )?;

    Ok(Self{
        dig_t1: u16::from_le_bytes([buffer[0], buffer[1]]),
        dig_t2: i16::from_le_bytes([buffer[2], buffer[3]]),
        dig_t3: i16::from_le_bytes([buffer[4], buffer[5]]),

        dig_p1: u16::from_le_bytes([buffer[6], buffer[7]]),
        // ...
        dig_p9: i16::from_le_bytes([buffer[22], buffer[23]]),
    })
}
\end{lstlisting}

\subsection*{Konfiguracja}
Konfiguracja modułu BMP280 odbywa się poprzez zapis 8-bitowych
wartości do rejestrów \texttt{CTRL\_MEAS} i~\texttt{CONFIG}. Wartości
te pozyskiwane są z~typu implementującego cechę \texttt{Config},
widoczną na listingu~\ref{lst:bmp280_config}. Cecha ta definiuje
skojarzone typy odpowiedzialne za poszczególne ustawienia modułu:
\begin{itemize}
    \item \texttt{TempOversampling} -- określa stopień nadpróbkowania
    pomiaru temperatury,
    \item \texttt{PressOversampling} -- określa stopień nadpróbkowania
    pomiaru ciśnienia,
    \item \texttt{PowerMode} -- określa tryb pracy urządzenia,
    \item \texttt{Standby} -- określa odstęp czasowy pomiędzy pomiarami,
    \item \texttt{IIR} -- określa stałą filtra IIR.
\end{itemize}

Każdy z~tych typów ograniczony jest do odpowiedniej cechy udostępniającej
stałą \texttt{BITS}, reprezentującą wartość bitową zgodną z~notą
katalogową. Wymagane wartości rejestrów przechowywane są w~stałych
\texttt{CTRL\_MEAS\_VALUE} i~\texttt{CONFIG\_VALUE}, których wartości
powstają przez przesunięcie bitowe stałych \texttt{BITS}
na odpowiednie pozycje i~złożenie ich operatorem alternatywy bitowej.
Rejestr \texttt{CONFIG\_VALUE} dodatkowo określa protokół komunikacyjny
urządzenia -- pole \texttt{EN\_SPI} ustawione jest na~0, co wymusza
użycie protokołu \isqrc.

\begin{lstlisting}[language=Rust, style=colouredRust, 
    caption={Cecha \texttt{Config} używana do konfiguracji BMP280},
    label={lst:bmp280_config}
]
pub trait Config{
    type TempOversampling: Oversampling;
    type PressOversampling: Oversampling;
    type PowerMode: PowerMode;
    type Standby: Standby;
    type IIR: IIRConstant;

    const CTRL_MEAS_VALUE: u8 = 
        Self::TempOversampling::BITS << TEMP_OVERSAMPLING_OFFSET
        | Self::PressOversampling::BITS << PRESS_OVERSAMPLING_OFFSET
        | Self::PowerMode::BITS << POWER_MODE_OFFSET;

    const CONFIG_VALUE: u8 =
        Self::Standby::BITS << STANDBY_OFFSET
        | Self::IIR::BITS << IIR_CONSTANT_OFFSET
        // Disable SPI mode
        | 0b0 << EN_SPI_OFFSET;
}

pub trait Oversampling{ const BITS: u8; }
pub trait PowerMode{ const BITS: u8; }
pub trait Standby{ const BITS: u8; }
pub trait IIRConstant{ const BITS: u8; }
\end{lstlisting}

Przykładową implementację cechy \texttt{Config} oraz cech
odpowiedzialnych za poszczególne wartości bitowe przedstawiono
na listingu~\ref{lst:bmp280_config_examples}. Definicje pomocniczych
cech \texttt{Oversampling}, \texttt{PowerMode}, \texttt{Standby}
i~\texttt{IIRConstant} zostały pominięte, ponieważ zawierają się 
w~listingu~\ref{lst:bmp280_config}.

\begin{lstlisting}[language=Rust, style=colouredRust, 
    caption={Przykładowa implementacja cech konfiguracyjnych modułu BMP280},
    label={lst:bmp280_config_examples}
]
pub struct DefaultConfig;

impl Config for DefaultConfig{
    type TempOversampling = OversamplingX1;
    type PressOversampling = OversamplingX2;
    type PowerMode = PowerModeNormal;

    type Standby = Standby0_5ms;
    type IIR = IIROff;
}

pub struct OversamplingX1;
pub struct OversamplingX2;
impl Oversampling for OversamplingX1{  const BITS: u8 = 0b001; }
impl Oversampling for OversamplingX2{  const BITS: u8 = 0b010; }

pub struct PowerModeNormal;
impl PowerMode for PowerModeNormal{ const BITS: u8 = 0b11; }

pub struct Standby0_5ms;
impl Standby for Standby0_5ms{   const BITS: u8 = 0b000; }

pub struct IIROff;
impl IIRConstant for IIROff{ const BITS: u8 = 0b000; }
\end{lstlisting}